% ============================================================================
% Appendix: Phase 2 N7 — Thick-Junction Internal-Core Node-Well Test
% ============================================================================
% Status: [Dc|model] derivation/bounding attempt within declared model class
% Scope: Test whether regularized thick-junction core energy at scale δ
%        generates a non-geometric attractive sector in V(q) that could
%        create a secondary minimum, beyond the accepted geometric baseline.
% Phase 2 / N7 WP2 deliverable
% Governing documents: PHASE2_NEXTSTEP_PLAN_V1.md,
%                      audit/N7_WP1_DONOR_NORMALIZATION.md
% ============================================================================

\chapter{Thick-Junction Internal-Core Test for the Node-Well Sector}
\label{app:P2_N7_core}

% ============================================================================
\section{Scope and Framing}
\label{sec:P2N7:scope}
% ============================================================================

This appendix tests the \textbf{N7 thick-junction / internal-core lane}
as a source of non-geometric attraction in the effective potential $V(q)$.
It builds on the accepted geometric single-well baseline
$V_{\mathrm{geom}}(q) = \tau L_{\mathrm{tot}}(q)$ \tagDc\
(Appendix~\ref{app:vq_chosen_path}) and follows the N7~WP1 donor
normalization and model boundary
(\texttt{audit/N7\_WP1\_DONOR\_NORMALIZATION.md}).

\textbf{What this appendix tests:}
Does the internal energy of a regularized Y-junction core (at scale
$\delta$) produce a $q$-dependent contribution $V_{\mathrm{core}}(q)$
that can oppose the geometric restoring force and create a secondary
minimum in the combined potential?

\textbf{What this appendix does not claim:}
\begin{itemize}
    \item It does not claim full closure of $V(q)$.
    \item It does not determine the core profile $f(q/\delta)$
          from first principles.
    \item It does not resolve the $L_0/\delta$ ratio.
    \item A no-go is an allowed and informative outcome.
\end{itemize}

% ============================================================================
\section{Donor Discipline}
\label{sec:P2N7:donors}
% ============================================================================

\textbf{Allowed donor bundle} (from N7~WP1 normalization):
\begin{itemize}
    \item Separable core ansatz structure (D-N7-1) \tagDc\ structural
    \item $V_{\mathrm{geom}}(q) = \tau L_{\mathrm{tot}}(q)$ (D-N7-2)
          \tagDc
    \item $\sigma = 8.82\MeV/\fm^2$ \tagDc
    \item $\delta = \hbar/(2m_p c) \approx 0.105\fm$ \tagI
    \item $L_0 \approx 1.0\fm$ \tagI
    \item $E_0 = \sigma L_0^2$ \tagDc
    \item WP2 Israel no-go (S-N7-3) \tagDc\ constraint
\end{itemize}

\textbf{Forbidden imports:}
\begin{itemize}
    \item $C = 100$ as independent evidence (F-N7-1)
    \item Phenomenological Gaussian/Lorentzian profiles as [Dc] (F-N7-2)
    \item Variant~3 fitted parameters ($V_0$, $q_*$, $w$)
    \item $V_B = 2\Delta m_{np}$ as derivation constraint
    \item $\tau_n \approx 878\second$ as input to $V(q)$
    \item N1 Israel thin-junction energy (bounded no-go)
\end{itemize}

\textbf{Anti-smuggling test} (from N7~WP1 \S8.3):
\emph{Delete all knowledge of $\tau_n$, $V_B \approx 2.6\MeV$, and the
Variant~3 Gaussian from the derivation. Does the derivation still stand?}

% ============================================================================
\section{Model-Class Setup}
\label{sec:P2N7:model}
% ============================================================================

\subsection{What ``Thick Junction'' Means}

In the thin-junction approximation (WP2,
Appendix~\ref{app:P2_WP2_Israel}), the Y-junction vertex is treated as
a distributional (zero-width) matching condition. The Israel conditions
apply at the vertex, and the result is $\Delta\theta \equiv 0$ with
arm-interior energy $\propto V_{\mathrm{geom}}$ (single-well).

In the \textbf{thick-junction model}, the vertex is regularized to a
finite-size core of characteristic radius $\delta \approx 0.105\fm$ \tagI.
Inside this core, the three worldsheets smoothly merge rather than
meeting at a mathematical point. The energy stored in this merging
region depends on the stress configuration, which changes when the
junction node is displaced by $q$.

\subsection{What ``Internal Core'' Means}

The internal core is the region $r_\perp < r_0$ around the junction
vertex, where $r_\perp$ is the transverse distance from the vertex and
$r_0 = L_0$ \tagI\ is the transverse core radius (identified with the
junction footprint). The core extends a distance $\sim \delta$ in the
$q$-direction (the collective coordinate of node displacement).

\subsection{Scales and Assumptions}

\begin{assumptionbox}[title=N7 Model Assumptions]
\begin{enumerate}
    \item \textbf{Separable core density} [P]:
          $\rho_{\mathrm{core}}(r_\perp, q) = \sigma \times
           g_\perp(r_\perp / r_0) \times f(q/\delta)$
    \item \textbf{Core scale} $\delta = \hbar/(2m_p c) \approx 0.105\fm$
          [I] (Compton anchor)
    \item \textbf{Transverse extent} $r_0 = L_0 \approx 1.0\fm$ [I]
          (junction footprint)
    \item \textbf{Transverse profile integral}
          $I_\perp = \int d^2\xi\, g_\perp(|\xi|) \approx \pi$
          for well-behaved profiles [Dc]
    \item \textbf{Adiabatic approximation}: fast modes relax at
          fixed $q$ [P]
\end{enumerate}
\end{assumptionbox}

\subsection{Core Energy: Structural Form}

From the separable ansatz (D-N7-1), integrating out transverse
directions:
\begin{equation}
    V_{\mathrm{core}}(q) = -E_0 \times f(q/\delta)
    \label{eq:N7_Vcore}
    \tagDc
\end{equation}
where
\begin{equation}
    E_0 = \sigma \times r_0^2 \times I_\perp
        = \sigma \times L_0^2 \times I_\perp
    \label{eq:N7_E0}
    \tagDc
\end{equation}

The energy scale $E_0$ is independent of $\delta$. With
$I_\perp \approx \pi$:
\begin{equation}
    E_0 \approx \pi \times 8.82 \times 1.0^2
        \approx 27.7\MeV
    \label{eq:N7_E0_value}
\end{equation}

If $I_\perp$ is absorbed into $f$ (as in the donor convention where
$f(0) = 1$ already includes the profile normalization):
\begin{equation}
    E_0 = \sigma L_0^2 \approx 8.82\MeV
    \label{eq:N7_E0_convention}
\end{equation}

\textbf{Key point:} The energy scale $E_0 \sim O(10\MeV)$ is set by
$\sigma$ and $L_0$ alone. It is of the right magnitude to compete
with $V_{\mathrm{geom}}$, which varies by $\sim \tau R \sim O(10\MeV)$
over the displacement range. The energy scale is \tagDc.

The profile $f(q/\delta)$ is \textbf{not determined} by the separable
ansatz. It encodes the $q$-dependence of the core energy and is the
central open quantity.

% ============================================================================
\section{Core Reduction: Sign and Monotonicity Analysis}
\label{sec:P2N7:reduction}
% ============================================================================

\subsection{Physical Argument for Profile Shape}

The core energy $V_{\mathrm{core}}(q)$ depends on the stress
configuration inside the regularized junction vertex. Two contributions
determine the core energy at displacement $q$:

\begin{enumerate}
    \item \textbf{Binding energy} (negative): energy saved by the
          three worldsheets merging in the core region rather than
          remaining separate. This is the dominant attractive contribution.
          At the Steiner configuration ($q = 0$), all three arms converge
          symmetrically in the core, producing maximal binding.
          As $q$ increases, the angular asymmetry reduces the efficiency
          of the merging, and the binding energy weakens.

    \item \textbf{Strain energy} (positive): energy cost of
          accommodating the stress concentration at the junction vertex.
          For an elastic medium, the strain energy density scales as
          $\sigma^2/Y$ (where $Y$ is the elastic modulus). At $q = 0$
          (Steiner), the $\Zthree$-symmetric stress distribution
          minimizes the total strain energy (by the convexity of
          $\sigma^2$: equal stress distribution minimizes $\sum \sigma_i^2$
          at fixed $\sum \sigma_i$). Asymmetric configurations
          ($q > 0$) have higher strain energy.
\end{enumerate}

Both contributions favor $q = 0$:
\begin{itemize}
    \item Binding energy is most negative at $q = 0$ (maximal overlap).
    \item Strain energy is least positive at $q = 0$ (symmetric stress).
\end{itemize}

\begin{derivationbox}[title=Physical Sign Determination]
Within the elastic-medium model of the regularized core, the core
energy $V_{\mathrm{core}}(q)$ satisfies:
\begin{equation}
    V_{\mathrm{core}}(0) \leq V_{\mathrm{core}}(q)
    \quad \text{for all } q > 0
    \label{eq:N7_Vcore_min_at_zero}
\end{equation}
The Steiner configuration is a \textbf{minimum} of the core energy.
\tagDcI

\textbf{Consequence for the profile:} In the parameterization
$V_{\mathrm{core}}(q) = -E_0 f(q/\delta)$ with $E_0 > 0$, this
requires:
\begin{equation}
    f(0) \geq f(q/\delta) \quad \text{for all } q > 0
    \label{eq:N7_f_peaked_at_zero}
\end{equation}
That is, $f$ is peaked at $q = 0$ and non-increasing for $q > 0$.
\end{derivationbox}

\textbf{Epistemic status of the sign argument:} This uses
(i)~the separable ansatz [P], (ii)~the elastic-medium analogy for
the brane [P], and (iii)~the convexity of $\sigma^2$ for strain energy
[Der]. The argument is physically motivated but model-dependent.
\tagDcI

\subsection{Thin-Junction Limit Recovery}

The N7 model must reduce to the WP2 no-go in the $\delta \to 0$
limit. For any profile $f$ with $\int f(q/\delta)\, dq = \delta
\int f(u)\, du = O(\delta)$:
\begin{equation}
    V_{\mathrm{core}}(q) = -E_0 f(q/\delta)
    \xrightarrow{\delta \to 0}
    -E_0 \times \delta \times c_f \times \delta_{\mathrm{Dirac}}(q)
    \label{eq:N7_thin_limit}
\end{equation}
where $c_f = \int f(u)\, du$ is a dimensionless constant. The core
contribution becomes a contact interaction at $q = 0$ with vanishing
total weight ($\propto \delta \to 0$). It does not produce a secondary
minimum, consistent with the WP2 no-go. \tagDc

% ============================================================================
\section{Monotone Profile No-Go}
\label{sec:P2N7:nogo}
% ============================================================================

This is the central structural result of this appendix.

\subsection{Statement}

\begin{theorem}[Monotone Profile No-Go]
\label{thm:monotone_nogo}
Let $V_{\mathrm{geom}}(q) = \tau L_{\mathrm{tot}}(q)$ be the geometric
baseline (Appendix~\ref{app:vq_chosen_path}), and let
$V_{\mathrm{core}}(q) = -E_0 f(q/\delta)$ where $E_0 > 0$ and
$f$ is any function satisfying:
\begin{enumerate}
    \item[(i)] $f(0) = \max_{q \geq 0} f(q/\delta)$
          \quad (peaked at $q = 0$)
    \item[(ii)] $f'(u) \leq 0$ for $u > 0$
          \quad (non-increasing for $q > 0$)
\end{enumerate}
Then the combined potential
$V(q) = V_{\mathrm{geom}}(q) + V_{\mathrm{core}}(q)$ has a unique
minimum at $q = 0$. No secondary minimum exists for $q > 0$.
\end{theorem}

\begin{proof}
Differentiate the combined potential:
\begin{equation}
    V'(q) = V_{\mathrm{geom}}'(q) + V_{\mathrm{core}}'(q)
    = \tau L_{\mathrm{tot}}'(q) - \frac{E_0}{\delta}\, f'\!\left(\frac{q}{\delta}\right)
    \label{eq:N7_Vprime}
\end{equation}

\textbf{Term 1:} $V_{\mathrm{geom}}'(q) = \tau L_{\mathrm{tot}}'(q) > 0$
for all $q \in (0, R)$.

This was established in Appendix~\ref{app:vq_chosen_path}: the Steiner
point is the unique minimum of $L_{\mathrm{tot}}(q)$, and
$L_{\mathrm{tot}}$ is strictly increasing for $q > 0$
(Theorem~\ref{thm:steiner_min}).

Explicitly, for the in-brane path:
\begin{equation}
    L_{\mathrm{tot}}'(q) = -1 + \frac{R + 2q}{\sqrt{R^2 + Rq + q^2}}
    \label{eq:N7_Ltot_prime}
\end{equation}
At $q = 0$: $L_{\mathrm{tot}}'(0) = 0$. For $q > 0$:
$(R + 2q)^2 = R^2 + 4Rq + 4q^2 > R^2 + Rq + q^2$ (since
$3Rq + 3q^2 > 0$), so $(R+2q)/\sqrt{R^2+Rq+q^2} > 1$, hence
$L_{\mathrm{tot}}'(q) > 0$. \tagDc

\textbf{Term 2:} $V_{\mathrm{core}}'(q) = -(E_0/\delta) f'(q/\delta)$.
By hypothesis (ii), $f'(q/\delta) \leq 0$ for $q > 0$, so
$V_{\mathrm{core}}'(q) \geq 0$.

\textbf{Combining:} For $q > 0$:
\begin{equation}
    V'(q) = \underbrace{V_{\mathrm{geom}}'(q)}_{> 0}
           + \underbrace{V_{\mathrm{core}}'(q)}_{\geq 0}
           > 0
    \label{eq:N7_Vprime_positive}
\end{equation}

Since $V'(q) > 0$ for all $q > 0$, the potential $V(q)$ is strictly
increasing on $(0, R)$. Its unique minimum on $[0, R)$ is at $q = 0$.
No secondary minimum exists.
\end{proof}

\subsection{What This No-Go Covers}

The monotone profile class includes all ``natural'' core profiles:
\begin{center}
\begin{tabular}{lll}
\toprule
\textbf{Profile} & \textbf{$f(u)$} & \textbf{Covered?} \\
\midrule
Gaussian & $e^{-u^2}$ & YES \\
Lorentzian & $1/(1+u^2)$ & YES \\
Squared Lorentzian & $1/(1+u^2)^2$ & YES \\
Exponential & $e^{-|u|}$ & YES \\
Disk (compact) & $\Theta(1-|u|)$ & YES \\
Any $f$ with $f(0) = \max$ and $f' \leq 0$ & --- & YES \\
\bottomrule
\end{tabular}
\end{center}

\textbf{All three donor mechanisms (A1, A2, A3) from the junction-core
execution report fall in this class.} They are all monotonically
decaying from $q = 0$ and therefore cannot produce metastability.

\subsection{What This No-Go Does Not Cover}

The no-go applies only to profiles satisfying conditions (i) and (ii).
Profiles that are \textbf{not} peaked at $q = 0$ --- i.e., profiles with
$f(q_*/\delta) > f(0)$ for some $q_* > 0$ --- are not covered. Such
profiles would represent core energy that is \emph{more attractive} at
$q = q_*$ than at the Steiner equilibrium.

\begin{edcopenbox}[title=Escape Route: Non-Monotone Core Energy]
Metastability from the thick-junction core would require a profile $f$
with a local maximum at some $q_* > 0$:
\begin{equation}
    f(q_*/\delta) > f(0)
    \label{eq:N7_escape}
\end{equation}

\textbf{Physical meaning:} The core energy is \emph{lower} (more
negative) when the junction node is displaced to $q = q_*$ than when
it sits at the Steiner equilibrium. This would mean the Steiner
configuration is \emph{not} the core energy minimum.

\textbf{Physical plausibility:} The sign argument in
\S\ref{sec:P2N7:reduction} provides a physical case against this:
both binding energy (maximal overlap at $q=0$) and strain energy
(minimized by $\Zthree$-symmetric stress at $q=0$) favor $q = 0$ as
the core energy minimum. A non-monotone profile would require a
mechanism that makes the core energy \emph{decrease} with increasing
asymmetry, which has no obvious elastic-medium analogue.

\textbf{Possible mechanism:} A qualitative change in core structure
at some critical displacement (e.g., a $\Zthree \to \Ztwo$ symmetry
transition inside the core that releases stress energy). This would
require solving the internal core dynamics, which is beyond the
separable ansatz framework.

\textbf{Status:} [OPEN]. Not ruled out, but physically unmotivated
within the natural model class.
\tagOpen
\end{edcopenbox}

% ============================================================================
\section{Depth and Shape Analysis}
\label{sec:P2N7:depth}
% ============================================================================

\subsection{Core Contribution Near $q = 0$}

For a monotone profile with $f(0) = 1$ and $f''(0) < 0$ (concave at
origin):
\begin{equation}
    V_{\mathrm{core}}(q) \approx -E_0
    \left(1 + \frac{1}{2} f''(0) \frac{q^2}{\delta^2}\right)
    \label{eq:N7_Vcore_expand}
\end{equation}

The core contribution to the curvature at $q = 0$:
\begin{equation}
    V_{\mathrm{core}}''(0) = -\frac{E_0}{\delta^2}\, f''(0) > 0
    \label{eq:N7_Vcore_curvature}
\end{equation}
(positive because $f''(0) < 0$ for peaked profiles).

This \textbf{adds} to the geometric curvature:
\begin{equation}
    V''(0) = V_{\mathrm{geom}}''(0) + V_{\mathrm{core}}''(0)
    = \frac{3\tau}{2R} + \frac{E_0 |f''(0)|}{\delta^2}
    \label{eq:N7_V_curvature_total}
\end{equation}

The core makes the well at $q = 0$ \textbf{steeper}, not shallower. It
\emph{reinforces} the geometric restoring force rather than opposing it.

\subsection{Core Energy Width}

The core contribution decays on scale $\delta \approx 0.105\fm$. For
the Gaussian profile: $V_{\mathrm{core}}$ drops to $1/e$ of its maximum
at $q = \delta$. Compared to the geometric scale $R \sim 1\fm$, the
core attraction is concentrated in a narrow region $q \lesssim \delta
\ll R$ around the Steiner minimum.

\subsection{Can the Core Flatten the Geometric Rise?}

Even though the core cannot create a secondary minimum (by the monotone
no-go), it could in principle flatten the rising part of $V(q)$ near
$q \approx \delta$, where the core attraction is decaying. The
combined curvature at $q = 0$ is increased, but the slope at
$q \approx \delta$ could be reduced relative to $V_{\mathrm{geom}}$
alone.

However, since $V'(q) > 0$ for all $q > 0$ (Theorem~\ref{thm:monotone_nogo}),
the potential is still strictly increasing. There is no inflection
point, no flattening to zero slope, and no reversal. The core merely
modifies the rate of increase at small $q$ without changing the
qualitative single-well structure.

% ============================================================================
\section{Comparison Against Geometric Baseline}
\label{sec:P2N7:comparison}
% ============================================================================

\subsection{Combined Potential}

\begin{equation}
    V(q) = \tau L_{\mathrm{tot}}(q) - E_0 f(q/\delta)
    \label{eq:N7_Vtotal}
\end{equation}

\textbf{At $q = 0$:}
\begin{equation}
    V(0) = 3\tau R - E_0
    \label{eq:N7_V0}
\end{equation}
The core deepens the minimum by $E_0 \sim 8.82\MeV$ (or
$\sim 27.7\MeV$ with the $I_\perp$ factor).

\textbf{At $q \gg \delta$:}
\begin{equation}
    V(q) \approx \tau L_{\mathrm{tot}}(q)
    \quad (f \to 0)
    \label{eq:N7_Vlarge}
\end{equation}
The core attraction has decayed; only the geometric restoring force
remains.

\subsection{Does the Core Term Create New Structure?}

\textbf{No.} The core term:
\begin{itemize}
    \item Does \textbf{not} oppose the geometric restoring force
          ($V_{\mathrm{core}}'(q) \geq 0$ for $q > 0$)
    \item Does \textbf{not} create a secondary minimum
    \item Does \textbf{not} create a barrier
    \item \textbf{Deepens} the primary well at $q = 0$
    \item \textbf{Steepens} the curvature at $q = 0$
\end{itemize}

\subsection{Comparison to Variant~3}

The Variant~3 phenomenological well (forbidden donor F-N7-2) used
$V_{\mathrm{node}} = -V_0 \exp(-(q-q_*)^2/2w^2)$ with $q_* = 2\fm$.
The critical difference is $q_* \neq 0$: the attraction is centered
\emph{away from} the Steiner point, which allows it to oppose the
geometric restoring force at $q \approx q_*$ and create a secondary
minimum.

The thick-junction core energy, by contrast, is centered at $q = 0$
(the junction vertex). It adds to, rather than opposes, the geometric
tendency toward $q = 0$.

\textbf{Structural diagnosis:} The Variant~3 node well works because it
has $V_{\mathrm{node}}'(q) < 0$ for $q < q_*$ (pulling toward $q_*$,
opposing $V_{\mathrm{geom}}'$ which pulls toward $q = 0$). The
thick-junction core has $V_{\mathrm{core}}'(q) > 0$ for $q > 0$
(pulling toward $q = 0$, in the same direction as $V_{\mathrm{geom}}'$).
This is why the core cannot substitute for Variant~3.

% ============================================================================
\section{Structural vs.\ Numeric Closure Firewall}
\label{sec:P2N7:firewall}
% ============================================================================

\textbf{What is structural derivation in this appendix:}
\begin{itemize}
    \item The separable ansatz form
          $V_{\mathrm{core}}(q) = -E_0 f(q/\delta)$ \tagDc\ structural
    \item The energy scale $E_0 = \sigma L_0^2 \sim O(10\MeV)$
          \tagDc\ (from dimensional closure)
    \item The transverse integral $I_\perp \approx \pi$ for standard
          profiles \tagDc
    \item The monotone profile no-go
          (Theorem~\ref{thm:monotone_nogo}) \tagDc
    \item The sign argument: both binding and strain energy favor
          $q = 0$ as core energy minimum \tagDcI
    \item The thin-junction limit recovery \tagDc
\end{itemize}

\textbf{What is only parameterized (not derived):}
\begin{itemize}
    \item The profile $f(q/\delta)$ itself --- shape remains [OPEN]
    \item Whether a non-monotone profile is physically realizable
          --- [OPEN]
    \item The separability of $\rho_{\mathrm{core}}(r_\perp, q)$
          --- assumed [P]
    \item The identification $r_0 = L_0$ --- [I]
\end{itemize}

\textbf{Why numeric barrier matching is not used:}
The value $V_B \approx 2.6\MeV$ from the calibrated $\tau_n$ formula
does not enter this analysis. No parameter in this appendix is adjusted
to reproduce $V_B$, $\tau_n$, or $\Delta m_{np}$. The conclusion (no
secondary minimum for monotone profiles) is \emph{independent} of these
target values.

\textbf{Why $C = 100$ reasoning is excluded:}
The numerical value $C = (L_0/\delta)^2$ is not cited as evidence.
Only the \emph{structural} scaling $E_0 = \sigma L_0^2$ is used,
and this is the $\delta$-independent reformulation from the delta
audit (N7~WP1 \S7.4).

% ============================================================================
\section{Outcome Classification}
\label{sec:P2N7:outcome}
% ============================================================================

\begin{resultbox}[title=N7 WP2 Outcome: Bounded Insufficiency]
\textbf{Classification:} Bounded insufficiency within the natural
model class.

\textbf{What was established:}
\begin{enumerate}
    \item \textbf{Energy scale [Dc]:} The thick-junction core energy
          scale $E_0 = \sigma L_0^2 \sim O(10\MeV)$ is sufficient to
          compete with $V_{\mathrm{geom}}$ in magnitude.
    \item \textbf{Monotone profile no-go [Dc]:} If the core profile $f$
          is peaked at $q = 0$ and monotonically non-increasing for
          $q > 0$ (the physically motivated class), then $V(q) =
          V_{\mathrm{geom}}(q) + V_{\mathrm{core}}(q)$ is a single-well
          potential. No secondary minimum can exist.
    \item \textbf{Sign result [Dc+I]:} Both binding energy (maximal
          overlap) and strain energy (symmetric stress) favor the Steiner
          configuration as the core energy minimum. The core contribution
          \emph{reinforces} the geometric restoring force.
    \item \textbf{Thin-junction limit [Dc]:} $V_{\mathrm{core}} \to 0$
          as $\delta \to 0$, recovering the WP2 no-go.
\end{enumerate}

\textbf{What this means:}
The thick-junction core, within the natural model class (monotone
profile, elastic core), cannot produce a secondary minimum. The core
energy deepens the Steiner well but does not create a competing
attractive region at $q > 0$.

\textbf{Why ``bounded insufficiency'' not ``no-go'':}
The full model class is not ruled out. Non-monotone profiles
($f(q_*/\delta) > f(0)$ for some $q_* > 0$) could in principle produce
metastability, but require a physical mechanism that makes the core
energy \emph{decrease} with displacement from Steiner. No such mechanism
has been identified within the elastic-medium framework.

\textbf{Scope of the insufficiency:}
\begin{itemize}
    \item Separable core ansatz with monotone profile
    \item In-brane node displacement (R3 path)
    \item Elastic-medium model for core stress
    \item Any value of $\sigma$, $\delta$, $L_0$, $R$
\end{itemize}

\textbf{Epistemic tag:} \tagDc\ (structural no-go within model class).
\end{resultbox}

\subsection{Impact on Candidate Landscape}

\begin{center}
\begin{tabular}{llll}
\toprule
\textbf{Candidate} & \textbf{Before N7 WP2} & \textbf{After N7 WP2}
& \textbf{Status} \\
\midrule
N1 Israel thin-junction & Bounded no-go & Unchanged & Dead end \\
N7 thick-junction (monotone) & Active & Bounded insufficiency
& Narrowed \\
N7 thick-junction (non-monotone) & Not tested & [OPEN] & Requires
new physics \\
N2 bulk backreaction & Backup & Promoted to primary & Active \\
\bottomrule
\end{tabular}
\end{center}

% ============================================================================
\section{Summary}
\label{sec:P2N7:summary}
% ============================================================================

\begin{resultbox}[title=N7 WP2 Summary]
\textbf{Question asked:} Can the thick-junction / internal-core
contribution generate a non-geometric attractive sector in $V(q)$?

\textbf{Answer:} \textbf{Not within the natural model class} (monotone
core profile peaked at Steiner equilibrium).

\textbf{Key results:}
\begin{center}
\begin{tabular}{lll}
\toprule
\textbf{Item} & \textbf{Result} & \textbf{Tag} \\
\midrule
Energy scale $E_0$ & $\sigma L_0^2 \sim O(10\MeV)$ --- sufficient
& [Dc] \\
Monotone no-go & $V'(q) > 0$ for $q > 0$ --- no secondary min
& [Dc] \\
Sign of core term & Core reinforces geometric restoring force
& [Dc+I] \\
Thin-junction limit & $V_{\mathrm{core}} \to 0$ as $\delta \to 0$
& [Dc] \\
Profile $f(q/\delta)$ & Undetermined (central open quantity)
& [OPEN] \\
Non-monotone escape & Physically unmotivated but not ruled out
& [OPEN] \\
\bottomrule
\end{tabular}
\end{center}

\textbf{N7 lane status after WP2:} Narrowed. The natural model class
(monotone profile) produces bounded insufficiency. The lane is not
dead --- non-monotone profiles or non-separable core physics could
evade the no-go --- but it requires new physical input not present
in the current donor base.

\textbf{Recommended next step:}
\begin{itemize}
    \item \textbf{Primary:} Pivot to N2 (bulk backreaction) as the next
          active lane. N2 has not been tested with junction-sourced
          perturbations and could produce attraction through a different
          mechanism (gravitational backreaction $V_{\mathrm{bulk}}(q)$
          is not constrained by the monotone no-go).
    \item \textbf{Secondary:} If a physical mechanism for non-monotone
          core energy is identified (e.g., $\Zthree \to \Ztwo$ core
          transition), revisit N7 with the specific model.
    \item \textbf{If both N2 and N7 fail:} The double-well postulate
          [P] is falsified within the $S_{\mathrm{EH}} + S_{\mathrm{NG}}$
          model class.
\end{itemize}

\textbf{Anti-smuggling compliance:}
\begin{itemize}
    \item[$\checkmark$] No forbidden donors imported
    \item[$\checkmark$] $C = 100$ not used as evidence
    \item[$\checkmark$] No profile shape imported as [Dc]
    \item[$\checkmark$] $\tau_n$, $V_B$, $\Delta m_{np}$ not used
    \item[$\checkmark$] Monotone no-go derived structurally, not fitted
    \item[$\checkmark$] Bounded insufficiency documented honestly
\end{itemize}
\end{resultbox}
